% TOP_PROBLEM: {"id": 20000980, "title": "A boundary-flux reduction for the mean square-grid sandpile identity", "category_id": 2, "category": {"id": 2, "name": "combinatorics", "display_name": "Combinatorics", "description": "Counting problems, graph theory, discrete structures.", "slug": "combinatorics", "order_index": 2, "created_at": "2026-07-31T15:26:25.671Z"}, "status": "partially_solved", "problem_number": "AIM-COMBINATORICS-0105", "set_id": 14, "statusQualification": "Specifies wired square boxes and height convention 0--3; asks existence as well as value of the source scaling-limit average."}
% FIRST_POSTED: 2026-10-01
% ENTRY_KIND: statement_only
% UnsolvedMath ID 20000980; source catalogue code AIM-COMBINATORICS-0105
% !TeX program = lualatex
% Definition and sources only; this file is not an LLM solution attempt.
\documentclass[11pt,a4paper]{article}
\usepackage[margin=25mm]{geometry}
\usepackage{fontspec}
\setmainfont{lmroman10-regular.otf}[BoldFont=lmroman10-bold.otf,
 ItalicFont=lmroman10-italic.otf,BoldItalicFont=lmroman10-bolditalic.otf]
\usepackage{amsmath,amssymb,mathtools}
\usepackage{unicode-math}
\setmathfont{latinmodern-math.otf}
\AtBeginDocument{\let\setminus\smallsetminus}
\usepackage{xurl,hyperref}
\hypersetup{colorlinks=true,urlcolor=blue}
\setcounter{secnumdepth}{0}
\setlength{\parindent}{0pt}
\setlength{\parskip}{5pt}
\setlength{\emergencystretch}{3em}
\begin{document}
\section*{A boundary-flux reduction for the mean square-grid sandpile identity}\label{20000980}
{\small UnsolvedMath ID 20000980 · AIM-COMBINATORICS-0105 · Combinatorics · AIM Workshop Problem Lists}
\subsection{Definitions and mathematical statement}
For $n\geq1$, let $V_n=[-n,n]^2\cap\mathbb Z^2$.
Form a finite wired graph by retaining square-lattice edges between vertices of $V_n$ and replacing every edge leaving $V_n$ by an edge to a single sink.
Thus every nonsink vertex has degree four, counting sink-edge multiplicity.
A stable configuration is a function $\eta:V_n\to\{0,1,2,3\}$.
Stabilization repeatedly topples vertices with at least four grains, sending grains along incident edges and discarding those reaching the sink.
The result is independent of legal toppling order.

The recurrent configurations are the recurrent states of the finite Markov chain that adds a grain at a uniformly chosen nonsink vertex and then stabilizes.
Under addition followed by stabilization, they form the sandpile group.
Let $e_n$ denote its identity, characterized by
$(e_n+\eta)^\circ=\eta$ for every recurrent $\eta$, where $\circ$ denotes stabilization.

Determine whether the averages
\[
 a_n=\frac{1}{(2n+1)^2}\sum_{v\in V_n}e_n(v)
\]
converge as $n\to\infty$, and compute their limit.
This is the mean-value question for the identity's square-lattice scaling limit with square domains and wired boundary.
Equivalently, any limiting rescaled height profile should have this spatial average, with normalized area measure on the unit square.
The configuration being averaged is the group identity, not a uniform random recurrent configuration; these are different statistics.
The specified exhaustion and boundary convention make the requested limiting number unambiguous.
\subsection{Short English statement}
Compute the limiting average height of the sandpile group identity on growing wired square grids.
\subsection{Sources}
\begin{itemize}
\item[S1] ULAM.AI and the UnsolvedMath Contributors, supplied problems.json, record 20000980 (AIM-COMBINATORICS-0105), AIM Workshop Problem Lists. Catalogue identity and source transcription; CC BY 4.0.
\url{https://huggingface.co/datasets/ulamai/UnsolvedMath}.
\item[S2] American Institute of Mathematics, Generalizations of chip-firing and the critical group, item 15. Original workshop question underlying AIM-COMBINATORICS-0105.
\url{https://aimath.org/pastworkshops/chipfiringproblems.pdf}.
\end{itemize}
\end{document}
